\section{Analytic inputs and the comparison framework}\label{sec:inputs}

The analytic argument has three sources with different roles.
Walsh's density-model and multiplicative-invariance framework
\cite[Sections 2--7]{walsh} supplies the architecture of the comparison.
He explicitly anticipates the sums-of-two-squares indicator application
in Section 1.2.4 and states on p.~2 that classical estimates could
replace the zero-free input in his fixed-shift divisor theorem.
Our completed fixed-shift proof below retains that input at large
Mellin heights; Section~\ref{sec:fixed-classical-boundary} separates
the classical range from the remaining estimate.
The zero-free theorem \cite[Theorem 3.1]{zerofree} supplies a fixed
half-plane for Dirichlet \(L\)-functions. The weighted-dilation theorem
\cite[Section 3]{weightedgraphs} transfers an independent-label operator
estimate to a moment on integers. We state the latter input with its
normalizations, since constancy on a mark fiber and the order of parameter
choices matter in this application.

\subsection{The zero-free input}

\begin{proposition}[Imported zero-free theorem]\label{prop:zf}
Every Dirichlet \(L\)-function, including \(\zeta\), is nonzero in
\[
                         \Re s>11/12.
\]
The pole at \(s=1\) of a principal character is allowed. Both conductor
and height are unrestricted.
\end{proposition}

This is the Dirichlet conclusion of
\cite[Theorem 3.1, printed p.~10]{zerofree}. The same theorem treats
finite-order Hecke characters over \(\Q(\sqrt{-3})\); that stronger
form is not needed here. Neither the boundary line nor the sharper
\(7/8\) assertion in Part II of that source is used. The input is
nonvanishing on the open half-plane, not an asserted numerical
threshold for the representation theorem. Estimates for individual
\(L\)-functions used below require the usual functional equation,
Euler product, and holomorphic logarithm on zero-free rectangles in
addition to this proposition.

\subsection{Prime groups and the two operators}

For this subsection only, write
\[
                \Lambda=\log x,\qquad T=\log\Lambda.
\]
Fix constants
\[
 0<\alpha<\beta<\gamma<\delta<0.47,\quad
 c_0,C_0,v_-,v_+>0,\quad 0<\vartheta<1,\quad \ell\ge1.
\]
There are disjoint collections \(\mathcal G_s,\mathcal G_b\) of
pairwise disjoint finite prime sets \(P_g\), with
\[
 c_0T\le|\mathcal G_s|,|\mathcal G_b|\le C_0T.
\]
They satisfy
\[
\begin{split}
 e^{\Lambda^\alpha}\le p\le e^{\Lambda^\beta}
       &\quad(g\in\mathcal G_s,\ p\in P_g),\\
 e^{\Lambda^\gamma}\le p\le e^{\Lambda^\delta}
       &\quad(g\in\mathcal G_b,\ p\in P_g),\\
 v_-\le V_g:=\sum_{p\in P_g}\frac1p\le v_+,
       &\qquad \mu_g(p)=\frac1{V_gp}.
\end{split}                                                    \tag{3.1}\label{eq:inputbanks}
\]
The sets need not be intervals. This freedom is needed when taking
identity residue classes and deleting primes dividing the target.

Fix a sufficiently large integer \(J\), and put \(M=J+\ell\).
A pattern \(\nu\) specifies subsets
\(C_g\subseteq\{1,\ldots,J\}\), with
\[
 C_g=\varnothing\quad(g\in\mathcal G_s),\qquad
 |C_g|\le m\quad(g\in\mathcal G_b),\qquad t_g=\ell+|C_g|.
\]
The slots in \(\{1,\ldots,J\}\setminus C_g\) are shared between the
two endpoint lists. The remaining \(t_g\) slots at each endpoint
are unshared. Form \(D\) from the shared labels and, in the
\(C_g\)-slots, fresh independent draws with laws \(\mu_g\).
Thus \(D\) has \(J\) prime factors from each group, counted with
multiplicity. Write \(D=D_sD_b\) according to group size.

Each pattern has a complex coefficient \(K_\nu\). It may depend on
the ordered big-group source labels, target labels, and auxiliary
labels used in \(D_b\), and on no small-group labels. Assume
\[
                 \sum_\nu\sup |K_\nu|\le\Lambda^{C_1},           \tag{3.2}\label{eq:inputmass}
\]
where \(m,C_1\) are fixed independently of \(J\). All operators below
average source and target slot permutations independently in each
applicable group. These averages are orthogonal projections.

The \emph{ideal operator} acts on
\[
 \mathcal H_{\mathrm{id}}=
 L^2\!\left(\prod_{g\in\mathcal G_b}P_g^M,
                    \bigotimes_{g\in\mathcal G_b}\mu_g^{\otimes M}\right).
\]
Numerical repetitions are permitted. Its row operation retains shared
labels, samples the unshared target labels and auxiliary labels
independently, and multiplies by
\[
                         K_\nu D_b^{\,i\zeta}e(\Theta D_b),
             \qquad e(t)=\exp(2\pi it).
\]
Sum over patterns and compose with the permutation projections on
both sides; denote the result by \(\mathcal T(\Theta)\).

A \emph{physical state} consists of an integer \(n\) and an ordered
\(M\)-tuple in every group, with distinct labels within each group,
all dividing \(n\). Give every such state at \(n\) mass
\(\prod_gV_g^{-M}\), and use counting measure on the integer position.
This defines \(\mathcal H_{\mathrm{ph}}\). Put
\[
 \omega_g(n)=\#\{p\in P_g:p\mid n\},\quad
 \omega(n)=\sum_g\omega_g(n),\quad
 s=|\mathcal G_s|+|\mathcal G_b|.
\]
For a fixed integer \(k\) with \(0<|k|\le\Lambda^{C_2}\), the
\emph{physical operator} \(A\) moves from \(n\) to \(n'=n+kD\).
It sums over the physical target lists with the prescribed shared
labels, using row normalization \(\prod_gV_g^{-t_g}\); auxiliary
labels retain their independent probability laws. Every auxiliary
label is forbidden from both endpoint lists, although auxiliary
labels may equal each other. Multiply the row by
\[
 K_\nu D^{\,i\zeta}
       \vartheta^{(\omega(n)-Ms)/2}
       \vartheta^{(\omega(n')-Ms)/2},                            \tag{3.3}\label{eq:inputphysical}
\]
sum over patterns and apply all endpoint permutation projections.
The adjoint reverses the displacement and conjugates the coefficient
with the endpoint roles interchanged. The joint endpoint
normalization in each group is \(V_g^{-M-t_g}\) in either direction.

\begin{proposition}[Imported absolute bounds]\label{prop:graphabsolute}
Replace every elementary transition coefficient of \(A\) by its
absolute value before summing patterns and averaging permutations.
The resulting operator and its adjoint have row sums at most
\(\Lambda^{C_{\mathrm{abs}}}\). The exponent may depend on \(m,C_1\)
and the fixed group data, but is independent of \(J\).
The same assertion holds with integer position replaced by independent
uniform residues at all group primes.
\end{proposition}

This is \cite[Lemma 3.4]{weightedgraphs}. The independence from \(J\)
is part of the assertion. To see its source, once the shared labels
are fixed, the new target slots in a group cost at most
\[
       \sup_{z\ge0}V_g^{-t_g}(z+t_g)^{t_g}\vartheta^{z/2}.
\]
Here \(t_g\le\ell+m\). This is a fixed constant per group and its
product over \(O(T)\) groups is a fixed power of \(\Lambda\).
Reversing the joint endpoint measure gives the column bound.

\subsection{Root moments, residue replacement, and endpoints}

\begin{proposition}[Imported local transference]\label{prop:rootmoment}
For every fixed \(E>0\), there is \(E_{\mathrm{id}}>0\), depending
only on \(E\) and the fixed group data, such that the following holds.
Suppose
\[
                   \sup_{\Theta\in\mathbb R}
                        \|\mathcal T(\Theta)\|_{2\to2}
                              \le\Lambda^{-E_{\mathrm{id}}}.
                                                                    \tag{3.4}\label{eq:inputideal}
\]
After \(m,C_1\) have been fixed, take a sufficiently large fixed \(J\).
Put \(R=\lceil\Lambda^{0.52}\rceil\). For every fixed \(\eta>0\)
and \(I=[X,2X)\), with \(x^\eta\le X\le x^{1/\eta}\),
\[
 \frac1{|I\cap\Z|}\sum_{n\in I\cap\Z}
       \langle u_n,(AA^*)^Ru_n\rangle
                         \le\Lambda^{-2RE}.                    \tag{3.5}\label{eq:inputroot}
\]
Here \(u_n\) is one on every physical mark list at position \(n\)
and zero elsewhere. The threshold may depend on all the fixed
parameters, including \(J,C_2,\eta\). No bound on \(\zeta\) is
required beyond \eqref{eq:inputideal}.
\end{proposition}

This is \cite[Theorem 3.5]{weightedgraphs}. It imposes return of the
integer position. The final mark list need not equal the initial
list, and \(u_n\) is not normalized. Substituting a trace with
identical initial and final mark lists would change the statement.

The arithmetic replacement inside that theorem has the following
precise form, which also explains its scale limitation.
For a fixed compatible path with \(r=2R\) edges, let \(h_i\) be
the integer offsets and \(\mathcal A_i\) its active prime set at
visit \(i\). Set \(\vartheta_i=\vartheta\) at internal visits and
\(\vartheta_0=\vartheta_r=\sqrt{\vartheta}\).
The position-dependent factor is
\[
 W(n)=\prod_{i=0}^{r}
       \left(\prod_{p\in\mathcal A_i}1_{p\mid n+h_i}\right)
       \vartheta_i^{\,\#\{p\notin\mathcal A_i:
                     \ p\text{ is a group prime},\ p\mid n+h_i\}}.
                                                                    \tag{3.6}\label{eq:inputpath}
\]

\begin{proposition}[Imported residue replacement]\label{prop:graphcrt}
For every fixed \(B>0\), the average of \eqref{eq:inputpath} on
\(I\cap\Z\) equals its expectation with independent uniform residues
at all group primes, multiplied by \(1+O(\exp(-BrT))\).
An incompatible collection of active congruences has zero weight
in both models.
\end{proposition}

This is \cite[Lemma 3.6]{weightedgraphs}. For clarity, its relative
form retains all deterministic damping belonging to active primes.
The product \(Q_{\mathrm{act}}\) of active primes satisfies
\[
              \log Q_{\mathrm{act}}
                   =O_J(rT\Lambda^\delta)=o(\Lambda).
\]
After those factors are extracted, a Bonferroni expansion of degree
\(O_B(rT)\) controls the remaining damping. Its independent
expectation is at least \(\exp(-O(rT))\); its total interval
counting error is \(\exp(-\eta\Lambda+o(\Lambda))\).
The inequality \(0.52+\delta<1\) makes that error negligible
even after division by the active congruence density. Together
with Proposition~\ref{prop:graphabsolute}, this permits summation
over paths with absolute coefficients. These observations describe
the imported replacement and do not assert replacement on intervals
shorter than its stated range.

\begin{corollary}[Imported endpoint consequence]\label{cor:graphendpoint}
Under Proposition~\ref{prop:rootmoment}, let \(f\in\mathcal H_{\mathrm{ph}}\)
be supported on positions \(I=[X,2X)\), constant on the chosen mark
list at each position, with value \(f_n\). If
\[
 \sup_{n\in I}|f_n|\le\exp(O(\sqrt{\Lambda})),\qquad
          \|f\|,\|g\|\le X^{1/2}\Lambda^{C_3},
\]
then
\[
                    |\langle f,Ag\rangle|
                        \ll X\Lambda^{-E+2C_3}.                 \tag{3.7}\label{eq:inputendpoint}
\]
Only \(f\) must be constant on mark lists.
\end{corollary}

This is \cite[Corollary 3.11]{weightedgraphs}. Its proof partitions
\(I\) into intervals on the displacement scale
\(\exp(O_J(T\Lambda^\delta))\), uses spectral H\"older for
\(AA^*\), and bounds the positive matrix
\((\langle u_n,(AA^*)^Ru_m\rangle)\) by its trace.
It makes no equal-fiber-mass assumption.

\subsection{What must be proved for the binary application}

The propositions above do not state a binary correlation theorem.
Walsh's theorem \cite[Theorem 1.1]{walsh} concerns
\(d_k(n)d_l(n+h)\) for positive \emph{integer} divisor orders and a
fixed nonzero shift \(h\). His Section 1.2.4 identifies other
indicator correlations as possible extensions. The membership
indicators and the reflected line in the present article require
separate arguments.

There are three interfaces. First, the one-variable Euler series
has a fractional branch \((s-1)^{-\rho_i}\). Its density model must
retain every character trivial on \(H_i\); an average scalar density
would lose genuine main terms. Second, the independent-label
comparison must be constructed for the fixed progression filters
and for prime groups in an identity residue class. The coefficient
bound \eqref{eq:inputmass} and the restriction to big labels must
hold before the final \(J\) is chosen. In particular, the source's
ideal-kernel construction for unrestricted interval prime sets is
not silently imported for these progression filters. Third, the
relation is \(n+(N-n)=N\), with the target growing on the summation
scale. Its comparison terms, endpoint norms, and localization
errors must be uniform in that target. We prove these interfaces
in the following sections; none is a consequence of the root
moment alone.

The parameter order is consequently part of the proof: choose the
required final accuracy, choose the ideal accuracy supplied by
Proposition~\ref{prop:rootmoment}, construct the signed family and
its fixed \(m,C_1\), and only then choose \(J\). The threshold in
\(x\) is taken last. The modulus and subgroup data remain fixed
throughout.
